\documentclass[10pt,a4paper]{amsart}
\usepackage[margin=19mm]{geometry}
\usepackage{amsmath,amssymb,mathtools}
\usepackage[scr=rsfs,cal=euler]{mathalfa}
\usepackage{xfrac}
\usepackage[unicode,hidelinks,bookmarks=false]{hyperref}
\newcommand{\ie}{i.e.\ }
\newcommand{\cf}{cf.\ }
\newcommand{\resp}{respectively\ }
\newcommand{\Subsection}{\subsection}
\providecommand{\nobreakdash}{\nobreak\hbox{-}}
\setlength{\emergencystretch}{2em}
\multlinegap0pt
\usepackage{bm}
\usepackage{bbm}
\usepackage{dsfont}
\usepackage{xspace}
\usepackage{cancel}
\usepackage{mathtools}
\usepackage{amsthm}
\makeatletter
\@ifundefined{Thm}{\newtheorem{Thm}{Thm}}{}
\@ifundefined{Conjecture}{\newtheorem{Conjecture}[Thm]{Conjecture}}{}
\@ifundefined{Prop}{\newtheorem{Prop}[Thm]{Proposition}}{}
\makeatother
\title[Moment generating functions in combinatorial optimization: B]{Moment generating functions in combinatorial optimization: Bipartite matching}
\author{Johan W\"astlund}
\date{Source: arXiv:2602.07563v1 (CC BY 4.0)}
\begin{document}
\maketitle

\noindent\emph{Source and licence.} Moment generating functions in combinatorial optimization: Bipartite matching. Version arXiv:2602.07563v1 (CC BY 4.0), distributed under the Creative Commons Attribution 4.0 International licence, as recorded in the arXiv licence metadata for this exact version and shown on the abstract page https://arxiv.org/abs/2602.07563v1. Full rights notice: licenses/arxiv-2602.07563-CC-BY-4.0.txt.\par\smallskip

\noindent\textit{Editorial scope.} Counted unit: the Prop whose source label is \texttt{P:zeroFree3} in the article (source0.tex lines 6244--6246, source label \texttt{P:zeroFree3}), delivered together with the article's own complete original proof of it (source0.tex lines 6248--6282, one \texttt{proof} environment of 35 lines). No mathematics is written, repaired or re-derived: statement and proof are the article's own text line for line. Required context retained uncounted: C:zeroFree (lines 6084--6086). Every definition, display and label of those lines is retained unchanged. Omitted from this package and not redistributed: 1--6083, 6087--6243, 6247--6247, 6283--6499. Nothing inside a retained excerpt is omitted or altered: every retained body is delivered exactly as the article's lines are written, so no omission receipt of any kind accompanies this package. Citations: the retained excerpts carry no citation command at all, so no bibliography entry of the article is transcribed and no reference is redistributed. Reference adaptation: none was needed and none was made; every reference inside the delivered unit resolves inside it, so no unresolved reference or citation is produced. Printed slips kept literal: no printed word, symbol, hypothesis or number of the counted statement or of its proof was altered. Layout and class: the article's own class is \texttt{article} with packages including babel, amsmath, amssymb, amsthm, tikz, pgfplots, xcolor, url; the wrapper is the lane's pinned \texttt{amsart} editorial wrapper at 10pt on A4, which restates the theorem-like environments the retained text needs and adds the article's own macro definitions that the retained text needs (none), with extra wrapper packages (bm, bbm, dsfont, xspace, cancel, mathtools). The delivered numbering is the unit's own. Assets: no retained span contains a figure environment, a graphics-inclusion command, a PDF-page inclusion, a table or a TikZ picture; no image file is copied into this package, the package declares no asset dependency and no image file is redistributed with it. Licence: version arXiv:2602.07563v1 of this article is distributed under the Creative Commons Attribution 4.0 International licence, as recorded in the arXiv licence metadata for this exact version and shown on the abstract page https://arxiv.org/abs/2602.07563v1. A full rights notice is delivered at licenses/arxiv-2602.07563-CC-BY-4.0.txt. Characters: every retained excerpt is pure ASCII, so the delivered engine, XeLaTeX with its default Latin Modern text font, reports no missing character.
\begin{Conjecture} \label{C:zeroFree}
The moment generating function $F_{k,m,n}(t)$ has no complex zeros in the open disk of radius $mn/k$ centered at the origin.
\end{Conjecture}

\begin{Prop} \label{P:zeroFree3}
The moment generating function $F_{3,m,n}(t)$ has no zero in the disk \[\left|t\right| \leq \frac{mn}{3}.\]
\end{Prop}

\begin{proof}
A direct computation shows that $F_{3, m, n}(t)$ can be written as
\[ F_{3,m,n}(t) = \frac{P(t)}{Q(t)},\] where the denominator is 
\begin{multline} 
Q(t) = (mn-3t)\,((m-1)n - 2t) \, (m(n-1)-2t) \\ 
((m-2)n-t) \, ((m-1)(n-1) - t) \, (m(n-2) - t),
\end{multline}
and the numerator is 
\begin{multline} P(t) = m^2n^2(m-1)(n-1)
-mn(4mn-3m-3n+2)t \\
+ (5mn-2m-2n-1)t^2 + 2t^3.
\end{multline}
Written with the same denominator $Q(t)$, the numerator of the Janson function $J_{3, m, n}(t)$ is
\begin{multline}
P_J(t) = (mn-2t)\cdot ((m-1)n-t)\cdot (m(n-1)-t) \\
= m^2n^2(m-1)(n-1)-mn(4mn-3m-3n+2)t \\ + (5mn-2m-2n)t^2 + 2t^3.
\end{multline} 
The two polynomials $P(t)$ and $P_J(t)$ are strikingly similar and differ by exactly $t^2$. Therefore on a circle $\left|t\right| = R$ the difference between them is
\[ \left| P(t) - P_J(t)\right| = R^2.\]
Moreover, provided that $R<mn/2$, it follows from the factorization of $P_J(t)$ that on the circle of radius $R$,
\[\left|P_J(t)\right| \geq P_J(R),\] 
since each of the three factors is minimized in absolute value when $t=R$.

If we take $R=mn/3$ (which is the radius relevant to Conjecture~\ref{C:zeroFree}), we find that $P_J(R)>R^2$, since
\begin{multline}  
P_J\left(\frac{mn}3\right)
= \left(mn-\frac{2mn}3\right)\cdot \left((m-1)n-\frac{mn}3\right) \cdot \left(m(n-1)-\frac{mn}3\right) \\
= \frac{m^2n^2}3\cdot\left(\frac{2m}{3}-1\right) \cdot \left(\frac{2n}3-1\right) \geq \frac{m^2n^2}3 = 3R^2.
\end{multline} 
Notice for the last inequality that in order for the optimization problem to make sense, both $m$ and $n$ have to be at least 3.

It follows from these estimates that on the circle $\left|t\right| = mn/3$, 
\[ \left| P(t) - P_J(t)\right| <  \left| P_J(t)\right|.\] 
Since the three zeros of $P_J(t)$ are $t=mn/2$, $t=(m-1)n$ and $t=m(n-1)$, all of them outside the disk $\left|t\right| \leq mn/3$, it follows from basic complex analysis that the zeros of $P(t)$ are also outside this disk. As $t$ goes around the boundary of the disk, both $P(t)$ and $P_J(t)$ will wind zero times around the origin.  
\end{proof}

\end{document}
